
\documentclass[11pt]{article}
\usepackage[T1]{fontenc}
\usepackage[utf8]{inputenc}
\usepackage[ngerman]{babel}
\usepackage{lmodern}
\usepackage{amsmath,amssymb,mathtools,array,booktabs}
\usepackage{geometry}
\usepackage{microtype}
\geometry{a4paper,margin=1.45cm}
\setlength{\parindent}{1.25em}
\setlength{\parskip}{0pt}
\makeatletter
\renewcommand\footnotesize{\@setfontsize\footnotesize{7.5}{8.5}\selectfont}
\makeatother
\allowdisplaybreaks
\setlength{\emergencystretch}{12em}
\sloppy
\hbadness=10000
\vbadness=10000
\hfuzz=2pt
\vfuzz=10pt
\raggedbottom
\newcommand{\ideal}[1]{\mathfrak{#1}}
\newcommand{\eqzero}[1]{\equiv 0\,(#1)}
\begin{document}
\setcounter{footnote}{34}

\subsection*{\S~9. Ausdehnung der Untersuchung auf Moduln. Anzahlgleichheit der Komponenten bei Zerlegungen in irreduzible Moduln.}

Wir zeigen jetzt, daß der Inhalt der drei ersten Paragraphen, der sich auf irreduzible, nicht auf primäre und Prim-Ideale bezieht, unter geringeren Voraussetzungen bestehen bleibt. Diese Paragraphen benutzen nämlich das kommutative Gesetz der Multiplikation nicht und beziehen sich nur auf die Eigenschaft der Ideale, Moduln zu sein, bleiben also erhalten für Moduln in bezug auf nicht-kommutative Bereiche, die jetzt zu definieren sind. Der Definition der Moduln ist ein Doppelbereich $(\Sigma,T)$ zugrunde zu legen von folgenden Eigenschaften:

$\Sigma$ ist ein abstrakt-definierter nicht-kommutativer Ring, d. h. $\Sigma$ ist ein System von Elementen $a,b,c,\ldots$, für die zwei Verknüpfungsarten definiert sind, die Ringaddition $(+)$ und die Ringmultiplikation $(\times)$, die den in \S~1 aufgestellten Gesetzen genügen, mit Ausnahme des kommutativen Gesetzes 4 der Ringmultiplikation.

$T$ ist ein System von Elementen $\alpha,\beta,\gamma,\ldots$, für das in Verbindung mit $\Sigma$ ebenfalls zwei Verknüpfungen definiert sind: die Addition, die aus je zwei Elementen $\alpha,\beta$ eindeutig ein drittes $\alpha+\beta$ erzeugt; die Multiplikation eines Elementes $\alpha$ aus $T$ mit einem Element $c$ aus $\Sigma$, die eindeutig ein Element $c\alpha$ aus $T$ erzeugt.\footnote{Es handelt sich also hier um eine ``rechtsseitige'' Multiplikation, einen ``rechtsseitigen'' Bereich $T$ und folglich um ``rechtsseitige'' Moduln und Ideale. Würde man für $T$ eine linksseitige Multiplikation $\alpha c$ zugrunde legen, so käme eine entsprechende Theorie der linksseitigen Moduln und Ideale; $M$ enthält neben $\alpha$ auch $\alpha c$. Das assoziative Gesetz wäre hier von der Form $(\gamma b)a=\gamma(b\times a)$.}

Für diese Verknüpfungen gelten die folgenden Gesetze:
\begin{enumerate}
\item Das assoziative Gesetz der Addition: $(\alpha+\beta)+\gamma=\alpha+(\beta+\gamma)$.
\item Das kommutative Gesetz der Addition: $\alpha+\beta=\beta+\alpha$.
\item Das Gesetz der unbeschränkten und eindeutigen Subtraktion: es gibt in $T$ ein und nur ein Element $\xi$, das die Gleichung $\alpha+\xi=\beta$ befriedigt (man bezeichnet $\xi=\beta-\alpha$).
\item Das assoziative Gesetz der Multiplikation: $a(b\gamma)=(a\times b)\gamma$.
\item Das distributive Gesetz:
\[
 (a+b)\gamma=a\gamma+b\gamma,\qquad c(\alpha+\beta)=c\alpha+c\beta.
\]
\end{enumerate}
Aus diesen Bedingungen folgt bekanntlich die Existenz des Nullelements, und die Gültigkeit des distributiven Gesetzes auch für die Verknüpfung von Subtraktion und Multiplikation:
\[
 (a-b)\gamma=a\gamma-b\gamma,\qquad c(\alpha-\beta)=c\alpha-c\beta,
\]
wo $(-)$ die Subtraktion in $\Sigma$ bedeutet. Enthält $\Sigma$ eine Einheit $\varepsilon$, so soll $\varepsilon\alpha=\alpha$ gelten für jedes Element $\alpha$ aus $T$.

Unter einem Modul $M$ in $(\Sigma,T)$ sei ein System von Elementen aus $T$ verstanden, das den beiden Bedingungen genügt:
\begin{enumerate}
\item $M$ enthält neben $\alpha$ auch $c\alpha$, wo $c$ ein beliebiges Element aus $\Sigma$ ist.
\item $M$ enthält neben $\alpha$ und $\beta$ auch die Differenz $\alpha-\beta$, also neben $\alpha$ auch $n\alpha$ für jede ganze Zahl $n$.\footnote{Die ganzen Zahlen sind wieder als abkürzende Zeichen, nicht als Ringelemente zu betrachten.}
\end{enumerate}
Nach dieser Definition bildet $T$ selbst einen Modul in $(\Sigma,T)$. Fallen insbesondere der Bereich $T$ und die dort festgelegten Verknüpfungen mit dem Bereich $\Sigma$ und den dort geltenden Verknüpfungen zusammen, so geht der Modul $M$ über in ein rechtsseitiges Ideal $M$ in $\Sigma$. Wird $\Sigma$ noch als kommutativ angenommen, so entsteht der gewöhnliche Idealbegriff, der sich somit als Spezialfall des Modulbegriffs ergibt.\footnote{Das einfachste Beispiel eines Moduls bildet der Modul aus ganzzahligen Linearformen: $\Sigma$ besteht hier aus allen ganzen rationalen Zahlen, $T$ aus allen ganzzahligen Linearformen. Ein etwas allgemeinerer Modul entsteht, wenn man in $\Sigma$ und $T$ ganze algebraische Zahlen statt der ganzrationalen Zahlen nimmt, oder etwa alle geraden Zahlen. Betrachtet man statt der Linearformen jeweils den Komplex aller Koeffizienten als ein Element, so sind die Verknüpfungen in $\Sigma$ und $T$ tatsächlich verschiedene. Ideale in nicht-kommutativen Ringbereichen aus Polynomen bilden den Gegenstand der gemeinsamen Arbeit Noether--Schmeidler. Von Idealen in weiteren speziellen nicht-kommutativen Bereichen handeln die Vorlesungen über die Zahlentheorie der Quaternionen von Hurwitz (Berlin, Springer 1919) und die dort zitierten Arbeiten von Du Pasquier.}

Für Moduln bleiben alle Definitionen des \S~1 erhalten. So bedeutet $\alpha\eqzero{M}$ bzw. $N\eqzero{M}$, daß $\alpha$ bzw. jedes Element von $N$ Element von $M$ ist; anders ausgedrückt, $\alpha$ bzw. $N$ ist durch $M$ teilbar. $M$ wird ein echter Teiler von $N$, wenn $M$ von $N$ verschiedene Elemente enthält; aus $N\eqzero{M}$, $M\eqzero{N}$ folgt $M=N$. Auch die Definition des größten gemeinsamen Teilers und des kleinsten gemeinsamen Vielfachen bleibt wörtlich erhalten. Enthält der Modul $M$ insbesondere eine endliche Anzahl von Elementen $\alpha_1,\ldots,\alpha_\rho$, derart, daß
\[
 M=(\alpha_1,
\ldots,\alpha_\rho),
\]
d. h.
\[
 \alpha=c_1\alpha_1+\cdots+c_\rho\alpha_\rho
       +n_1\alpha_1+\cdots+n_\rho\alpha_\rho
\]
wird für jedes $\alpha\eqzero{M}$, wobei die $c_i$ Größen aus $\Sigma$, die $n_i$ ganze Zahlen sind, so heißt $M$ ein endlicher Modul, $\alpha_1,\ldots,\alpha_\rho$ eine Modulbasis.

Wir legen nun im folgenden, analog wie in \S~1, nur solche Bereiche $(\Sigma,T)$ zugrunde, die die Endlichkeitsbedingung erfüllen: jeder Modul in $(\Sigma,T)$ ist ein endlicher, besitzt also eine Modulbasis.

Dann gilt für diesen Bereich $(\Sigma,T)$ Satz I von der endlichen Kette auch für Moduln, wie der dortige Beweis zeigt, und damit sind alle Voraussetzungen für die \S\S~2 und 3 erfüllt. Es läßt sich Definition I und Hilfssatz I über kürzeste und reduzierte Darstellung direkt übertragen; ebenso Satz II über die Darstellbarkeit jedes Moduls als kleinstes gemeinsames Vielfaches von endlich vielen irreduziblen Moduln, wobei Hilfssatz II zeigt, daß jede solche kürzeste Darstellung zugleich reduziert ist. Weiter bleibt Satz III erhalten, der die Reduzibilität eines Moduls durch Eigenschaften seines Komplements ausdrückt; und daraus folgt Hilfssatz III und schließlich Satz IV, der die Anzahlgleichheit der Komponenten bei zwei verschiedenen kürzesten Darstellungen eines Moduls als kleinstes gemeinsames Vielfaches von irreduziblen Moduln aussagt. Satz IV ergibt noch den Hilfssatz IV als Umkehrung von Hilfssatz I über reduzierte Darstellung.

\emph{Zusatz.} Dieselbe Überlegung zeigt, daß alle diese Sätze und Definitionen erhalten bleiben, wenn man für zweiseitige Bereiche $T$, d. h. solche, die sowohl rechts- wie linksseitig sind, unter Modul durchweg einen zweiseitigen Modul versteht, d. h. einen solchen, der neben $\alpha$ auch $c\alpha$ und $\alpha c$ enthält, neben $\alpha$ und $\beta$ auch $\alpha-\beta$.

Während all diese Sätze sich nur auf den Begriff der Teilbarkeit und des kleinsten gemeinsamen Vielfachen stützen, beruhen die weiteren Eindeutigkeitssätze wesentlich auf dem Produktbegriff und lassen deshalb eine direkte Übertragung nicht zu. So läßt sich die Definition des primären und Prim-Ideals nicht auf Moduln übertragen, da das Produkt zweier Größen aus $T$ nicht definiert ist. Die formal mögliche Übertragung auf nicht-kommutative Ringe verliert aber ihren Sinn, da hier die Existenz des zugehörigen Primideals sich nicht beweisen läßt\footnote{Es folgt nämlich aus $p_1^\rho\eqzero{\ideal Q}$, $p_2^\sigma\eqzero{\ideal Q}$ hier nicht $(p_1-p_2)^{\rho+\sigma}\eqzero{\ideal Q}$, und ebenso folgt aus $(ab)^\rho\eqzero{\ideal Q}$ nicht $a^\rho b^\rho\eqzero{\ideal Q}$. Dadurch läßt sich also $\ideal P$ weder als Ideal nachweisen, noch kommt ihm die Eigenschaft der Primideale zu. Für zweiseitige Ideale läßt $\ideal P$ sich zwar als Ideal, nicht aber als Primideal nachweisen.} und auch der Beweis, daß ein irreduzibles Ideal primär ist, versagt. Diese beiden Tatsachen bilden aber die Grundlage der folgenden Eindeutigkeitssätze. Dagegen lassen sich, wenn der nicht-kommutative Ring eine Einheit besitzt, die teilerfremden und teilerfremd-irreduziblen Ideale definieren, und die zum Beweise von Satz II benutzten Überlegungen lassen erkennen, daß jedes Ideal sich als kleinstes gemeinsames Vielfaches von endlich vielen paarweise teilerfremden und teilerfremd-irreduziblen Idealen darstellen läßt.\footnote{Für spezielle, ``vollständig reduzible'' Ideale lassen sich auch hier Eindeutigkeitssätze aufstellen; vgl. die gemeinsame Arbeit Noether--Schmeidler.}

Schließlich sei noch ein hinreichendes Kriterium dafür erwähnt, daß in $(\Sigma,T)$ die Endlichkeitsbedingung erfüllt ist: enthält $\Sigma$ eine Einheit und erfüllt die Endlichkeitsbedingung, und ist $T$ selbst ein endlicher Modul in $(\Sigma,T)$, so ist jeder Modul in $(\Sigma,T)$ ein endlicher.

Aus der Existenz der Einheit in $\Sigma$ folgt nämlich für
\[
 \ideal M=(f_1,
\ldots,f_e),\qquad
 f=\bar b_1f_1+\cdots+\bar b_ef_e+n_1f_1+\cdots+n_ef_e,
\]
wo $\bar b_i$ Elemente aus $\Sigma$, $n_i$ ganze Zahlen sind, auch eine Darstellung
\[
 f=b_1f_1+\cdots+b_ef_e,
\]
wo $b_i$ Elemente aus $\Sigma$ sind. Denn wegen $f_i=\varepsilon f_i$ wird $n_if_i=n_i\varepsilon f_i$, und da $n_i\varepsilon=(\varepsilon+\cdots+\varepsilon)$ zu $\Sigma$ gehört, wird $b_i=\bar b_i+n_i\varepsilon$ ein Element aus $\Sigma$. Die Voraussetzung über $T$ sagt nun aus, daß jedes Element $\alpha$ aus $T$ eine Darstellung zuläßt
\[
 \alpha=\bar a_1\tau_1+\cdots+\bar a_k\tau_k+n_1\tau_1+\cdots+n_k\tau_k,
\]
und folglich
\[
 \alpha=a_1\tau_1+\cdots+a_k\tau_k,
\]
wo die zweite Darstellung wegen $\varepsilon\tau_i=\tau_i$ sich wie oben aus der ersten ergibt.

Durchlaufen nun die $\alpha$ die Elemente eines Moduls $M$ aus $(\Sigma,T)$, so durchlaufen die Koeffizienten $a_k$ von $\tau_k$ ein Ideal $\ideal M_k$ aus $\Sigma$; nach dem Obigen wird also für jedes $a_k\eqzero{\ideal M_k}$ auch
\[
 a_k=b_1a_k^{(1)}+\cdots+b_ea_k^{(e)}.
\]
Bezeichnet $\alpha^{(i)}$ ein Element aus $M$, für das der Koeffizient von $\tau_k$ gleich $a_k^{(i)}$ wird, so wird $\alpha-b_1\alpha^{(1)}-\cdots-b_e\alpha^{(e)}$ ein zu $M$ gehöriges Element, das nur von $\tau_1,
\ldots,\tau_{k-1}$ abhängt. Auf die Gesamtheit dieser $\tau_k$ nicht mehr enthaltenden Elemente, die einen Modul $M'$ bilden, läßt sich das Verfahren wiederholen, so daß durch endlich oftmalige Wiederholung die Behauptung bewiesen ist.

\subsection*{\S~10. Spezialfall des Polynombereichs.}

1. Der zugrunde gelegte Ringbereich $\Sigma$ bestehe aus allen Polynomen von $x_1,
\ldots,x_n$ mit beliebigen komplexen Koeffizienten, für den nach dem Hilbertschen Theorem von der Modulbasis (Ann. 36) die Endlichkeitsbedingung erfüllt ist. Es handelt sich um den Zusammenhang unserer Sätze mit den bekannten Sätzen der Eliminations- und Modultheorie.

Dieser Zusammenhang wird hergestellt durch den folgenden Spezialfall eines bekannten Hilbertschen Satzes:\footnote{Über die vollen Invariantensysteme. Math. Ann. 42 (1893), \S~3, S. 313.}

Verschwindet $f$ für jedes endliche Wertsystem von $x_1,
\ldots,x_n$, das Nullstelle aller Polynome eines Primideals $\ideal P$ (Nullstelle von $\ideal P$) ist, so ist $f$ durch $\ideal P$ teilbar. M. a. W.: ein Primideal $\ideal P$ besteht aus der Gesamtheit der Polynome, die in seinen Nullstellen verschwinden.\footnote{Dieser Spezialfall läßt sich, wie Lasker [Math. Ann. 60 (1905), S. 607] gezeigt hat, im Fall homogener Formen auch direkt beweisen, und es folgt dann umgekehrt hieraus wieder der Hilbertsche Satz im homogenen und inhomogenen Fall, den wir so aussprechen können: Verschwindet ein Ideal $\ideal R$ in allen Nullstellen von $M$, so ist eine Potenz von $\ideal R$ durch $M$ teilbar. Dieser Satz, und ebenso der Spezialfall, gilt aber nur, wenn der Wertevorrat der $x$ algebraisch abgeschlossen ist, kann daher aus unsern Sätzen allein nicht folgen, sondern muß die Wurzelexistenz benutzen; etwa an der Stelle, daß ein Ideal, dessen Nullstellen nur aus $x_1=0,
\ldots,x_n=0$ bestehen, alle Potenzprodukte der $x$ von einer gewissen Dimension an enthält. Der übrige Beweis läßt sich unter Benutzung unserer Sätze gegenüber Lasker etwas vereinfachen. Lasker muß nämlich auch zum Nachweis der Zerlegung eines Ideals in größte primäre den Hilbertschen Satz heranziehen.}

Verschwindet somit ein Produkt $fg$ für alle Nullstellen von $\ideal P$, so verschwindet mindestens ein Faktor; die Nullstellen bilden ein irreduzibles algebraisches Gebilde. Legt man umgekehrt diese Definition des irreduziblen Gebildes zugrunde, so zeigt sich, daß die Gesamtheit der in einem irreduziblen Gebilde verschwindenden Polynome ein Primideal bilden; Primideal und irreduzibles Gebilde entsprechen sich somit eineindeutig. Wegen $\ideal P^\rho\eqzero{\ideal Q}$, $\ideal Q\eqzero{\ideal P}$ stimmen ferner die Nullstellen eines primären Ideals mit denen seines zugehörigen Primideals überein.\footnote{Diese Eigenschaft eines primären Ideals, ein irreduzibles Gebilde zu besitzen, nimmt Macaulay (vgl. Einleitung) zur Definition, während Lasker nur den Begriff der Mannigfaltigkeit eines Gebildes in die Definition aufnimmt, im übrigen abstrakt definiert. Die nur für $x_1=0,
\ldots,x_n=0$ verschwindenden primären Ideale nehmen bei Lasker eine Sonderstellung ein.} Die Darstellung eines Ideals als kleinstes gemeinsames Vielfaches von größten primären Idealen ergibt also eine Auflösung aller Nullstellen des Ideals in irreduzible Gebilde; und wie Lasker gezeigt hat, gilt auch das Umgekehrte. Der Eindeutigkeitsnachweis der zugehörigen Primideale entspricht also hier dem Fundamentalsatz der Eliminationstheorie von der eindeutigen Zerlegbarkeit eines algebraischen Gebildes in irreduzible Gebilde, kann für spezielle Polynombereiche, wo keine eindeutige Produktdarstellung der Polynome durch irreduzible Polynome des Bereichs und folglich auch keine Eliminationstheorie besteht, als Äquivalent dieses Satzes der Eliminationstheorie gelten.

Die den isolierten primären Idealen entsprechenden irreduziblen Gebilde sind genau die bei der ``Minimalresolvente''\footnote{Vgl. etwa J. König, Einleitung in die allgemeine Theorie der algebraischen Größen (Leipzig, Teubner, 1903), S. 235.} auftretenden; denn die Nullstellen jedes nicht-isolierten größten primären Ideals sind zugleich Nullstellen mindestens eines isolierten, nämlich eines solchen, dessen zugehöriges Primideal durch das des nicht-isolierten Ideals teilbar ist. Die Eindeutigkeit der isolierten primären Ideale ergibt also in den Exponenten neue invariante Multiplizitätszahlen. Auch die Eindeutigkeitssätze bei der Auflösung der primären Ideale in irreduzible können als Ergänzung der Eliminationstheorie im Sinn der Multiplizität angesehen werden.

Nach diesen Bemerkungen lassen sich die verschiedenen Zerlegungen in ihrem Verhalten zu den algebraischen Gebilden deuten. Den paarweise teilerfremden Idealen entsprechen solche Gebilde, die keine gemeinsame Nullstelle besitzen; bei den gegenseitig relativprimen Idealen ist kein irreduzibles Gebilde des einen Ideals zugleich gemeinsame Nullstelle; die größten primären Ideale verschwinden nur in irreduzibeln Gebilden, die alle voneinander verschieden sind; bei der Zerlegung in irreduzible Ideale können dieselben irreduziblen Gebilde auch mehrfach auftreten.

Bemerkt sei noch, daß an Stelle des allgemeinen Polynombereichs auch der Bereich aller homogenen Formen zugrunde gelegt werden kann, denn man überzeugt sich leicht, daß auch bei den dort geltenden Verknüpfungen - die Addition ist nur für Formen gleicher Dimensionen definiert - die allgemeinen Sätze erhalten bleiben.\footnote{Daß hier aber im Falle der Mehrdeutigkeit neben homogenen auch inhomogene Zerlegungen existieren können, zeigt das Beispiel $ (x^3,xy,y^3)=[(x^3,y),(y^3,x)]=[(xy,x^3,y^3,x+y^3),(xy,x^3,y^3,y+x^2)]$.}

Ein einfachstes Beispiel der vier verschiedenen Zerlegungen - für das die Formeln unten folgen - ist nach dem obigen etwa gegeben durch eine Gerade und zwei dazu windschiefe, sich schneidende Gerade, von denen eine einen vom Schnittpunkt verschiedenen Punkt in höherer Multiplizität enthält. Der Zerlegung in teilerfremd-irreduzible Ideale entspricht die Zerlegung in die Gerade und das dazu windschiefe Gebilde; dies Gebilde zerfällt bei der Zerlegung in relativprim-irreduzible Ideale in die beiden Geraden; der Zerlegung in größte primäre Ideale entspricht eine Ablösung des Punktes höherer Multiplizität, während die Zerlegung in irreduzible Ideale eine Auflösung dieses Punktes bedingt.

Nimmt man diesen Punkt zum Anfangspunkt, die durchlaufende Gerade zur $y$-Achse, die diese schneidende Gerade der $x$-Achse parallel, die dazu windschiefe Gerade der $z$-Achse parallel, so wird eine solche Konfiguration etwa dargestellt durch die folgenden irreduziblen Ideale:\footnote{Davon sind die drei ersten als Primideale irreduzibel; $\ideal B_4$, da es nur die Teiler $(x^2,y,z)$ und $(x,y,z)$ besitzt; $\ideal B_5$, da jeder Teiler das Polynom $xy$ enthält.}
\[
\begin{gathered}
 \ideal B_1=(x-1,y),\quad
 \ideal B_2=(y-1,z),\quad
 \ideal B_3=(x,z),\quad
 \ideal B_4=(x^3,y,z),\quad
 \ideal B_5=(x^2,y^2,z). 
\end{gathered}
\]
Die zugehörigen Primideale werden:
\[
 \ideal P_1=\ideal B_1,
 \quad \ideal P_2=\ideal B_2,
 \quad \ideal P_3=\ideal B_3,
 \quad \ideal P_4=\ideal P_5=(x,y,z).
\]
Die größten primären Ideale werden:
\[
 \ideal Q_1=\ideal B_1,
 \quad \ideal Q_2=\ideal B_2,
 \quad \ideal Q_3=\ideal B_3,
 \quad \ideal Q_4=[\ideal B_4,\ideal B_5]=(x^3,y^2,x^2y,z).
\]
Die relativprim-irreduziblen Ideale werden:
\[
 \ideal R_1=\ideal Q_1,
 \quad \ideal R_2=\ideal Q_2,
 \quad \ideal R_3=[\ideal Q_3,\ideal Q_4]=(x^3,x^2y,xy^2,z).
\]
Die teilerfremd-irreduziblen Ideale werden:
\[
 \ideal S_1=\ideal R_1,
 \quad
 \ideal S_2=[\ideal R_2,\ideal R_3]=((y-1)x^3,(y-1)x^2y,(y-1)xy^2,z).
\]
Daraus kommt das Gesamtideal:
\[
\begin{aligned}
 \ideal M&=[\ideal S_1,\ideal S_2]=\ideal S_1\ideal S_2 \\
 &=((x-1)(y-1)x^3,(y-1)x^2y,(y-1)xy^2,(x-1)z,y(y-1)x^3,yz),
\end{aligned}
\]
wobei
\[
 1=-(y-1)x^3+(y-1)(x^3-1)+y,
\]
\[
 (y-1)(x^3-1)+y\eqzero{\ideal S_1},
 \qquad -(y-1)x^3\eqzero{\ideal S_2}.
\]
Dabei sind die Ideale $\ideal B_1$, $\ideal B_2$, $\ideal B_3$ isolierte, also auch bei den Zerlegungen in irreduzible und größte primäre Ideale eindeutig bestimmt. Die Ideale $\ideal B_4$ und $\ideal B_5$ bzw. $\ideal Q_4$ sind nicht-isolierte; sie sind nicht eindeutig bestimmt, sondern etwa ersetzbar durch
\[
 \ideal D_4=(x^3,y+\lambda x^2,z),
 \qquad
 \ideal D_5=(x^2+\mu xy,y^2,z),
\]
ebenso ist $\ideal Q_4$ ersetzbar durch
\[
 \overline{\ideal Q}_4=(x^3,x^2y,y^2+\lambda xy,z).
\]

2. Ebenso wie der allgemeine (und der ganzzahlige) Polynombereich erfüllt auch jeder endliche Integritätsbereich aus Polynomen - wie das Hilbertsche Theorem von der Modulbasis zeigt - die Endlichkeitsbedingung,\footnote{Ist umgekehrt für einen Polynombereich die Endlichkeitsbedingung erfüllt, und läßt jedes Polynom mindestens eine Darstellung zu, wo die Multiplikatoren von geringerem Grad in $x$ sind, so ist der Bereich ein endlicher Integritätsbereich.} wobei die Koeffizienten einem beliebigen Körper zugewiesen werden können. Wir geben noch ein Beispiel für den Bereich aller geraden Polynome, als dem einfachsten Bereich, wo wegen $x^2y^2=(xy)^2$ keine eindeutige Produktdarstellung der Polynome durch irreduzible Polynome des Bereichs existiert. Es handelt sich um die gleiche Konfiguration wie in dem obigen Beispiel, die jetzt gegeben wird durch die irreduziblen Ideale:
\[
\begin{gathered}
 \ideal B_1=(x^2-1,xy,y^2,yz),\quad
 \ideal B_2=(y^2-1,xz,yz,z^2),\\
 \ideal B_3=(x^2,xy,xz,yz,z^2),\quad
 \ideal B_4=(x^4,x^3y,y^2,xz,yz,z^2),\\
 \ideal B_5=(x^2,y^2,xz,yz,z^2).
\end{gathered}
\]
Die zugehörigen Primideale werden:
\[
 \ideal P_1=\ideal B_1,
 \quad \ideal P_2=\ideal B_2,
 \quad \ideal P_3=\ideal B_3,
 \quad \ideal P_4=\ideal P_5=(x^2,xy,y^2,xz,yz,z^2).
\]
Daß $\ideal P_1$ Primideal wird, folgt daraus, daß jedes Polynom des Bereichs von der Form wird
\[
 f=\varphi(z^2)+xz\,\psi(z^2)\pmod{\ideal P_1}.
\]
Sei also
\[
 f_1=\varphi_1(z^2)+xz\,\psi_1(z^2),
 \qquad
 f_2=\varphi_2(z^2)+xz\,\psi_2(z^2),
\]
so folgt aus $f_1f_2\eqzero{\ideal P_1}$ das Bestehen der Gleichungen
\[
 \varphi_1\varphi_2+z^2\psi_1\psi_2=0,
 \qquad
 \varphi_2\psi_1+\varphi_1\psi_2=0,
\]
und daraus $f_1\eqzero{\ideal P_1}$ oder $f_2\eqzero{\ideal P_1}$. Genau so zeigt sich, daß $\ideal P_2$ Primideal wird; $\ideal P_3$ wird Primideal, da jedes Polynom des Bereichs modulo $\ideal P_3$ einem Polynom von $y^2$ kongruent wird; $\ideal P_4$ besteht aus allen Polynomen des Bereichs. Es werden also auch $\ideal B_1$, $\ideal B_2$ und $\ideal B_3$ als Primideale irreduzibel; $\ideal B_4$ und $\ideal B_5$ besitzen aber je nur den einzigen echten Teiler $\ideal P_4$, sind also notwendig irreduzibel.

Aus den irreduziblen Idealen ergeben sich die größten primären:
\[
\begin{gathered}
 \ideal Q_1=\ideal B_1,
 \quad \ideal Q_2=\ideal B_2,
 \quad \ideal Q_3=\ideal B_3,
 \quad
 \ideal Q_4=[\ideal B_4,\ideal B_5]=(x^4,x^3y,y^2,xz,yz,z^2),
\end{gathered}
\]
die relativprim-irreduziblen Ideale:
\[
 \ideal R_1=\ideal Q_1,
 \quad \ideal R_2=\ideal Q_2,
 \quad
 \ideal R_3=[\ideal Q_3,\ideal Q_4]=(x^4,x^3y,x^2y^2,xy^3,xz,yz,z^2),
\]
die teilerfremd-irreduziblen Ideale:
\[
 \ideal S_1=\ideal R_1,
 \quad \ideal S_2=[\ideal R_2,\ideal R_3]
 =((y^2-1)x^4,(y^2-1)x^3y,(y^2-1)x^2y^2,
 (y^2-1)xy^3,xz,yz,z^2).
\]
Es wird wie im ersten Beispiel:
\[
 [\ideal B_4,\ideal B_5]=[\ideal D_4,\ideal D_5],
 \qquad
 [\ideal Q_3,\ideal Q_4]=[\ideal Q_3,
\overline{\ideal Q}_4],
\]
wo
\[
 \ideal D_4=(x^4,xy+\lambda x^2,\ldots),\qquad
 \ideal D_5=(x^2+\mu xy,\ldots),\qquad \lambda\mu\ne1,
\]
und
\[
 \overline{\ideal Q}_4=(x^4,x^3y,y^2+\lambda xy,\ldots).
\]
Die übrigen irreduziblen bzw. größten primären Ideale $\ideal B_1$, $\ideal B_2$, $\ideal B_3$ sind als isolierte Ideale eindeutig bestimmt.

\subsection*{\S~11. Beispiele aus der Zahlentheorie und der Theorie der Differentialausdrücke.}

1. Der zugrunde gelegte Bereich $\Sigma$ bestehe aus allen geraden, ganzen rationalen Zahlen. $\Sigma$ läßt sich also eineindeutig dem Bereich aller ganzen rationalen Zahlen zuordnen, indem man jeder Zahl $2a$ aus $\Sigma$ die Zahl $a$ entsprechen läßt. Daraus folgt sofort, daß jedes Ideal in $\Sigma$ ein Hauptideal $(2a)$ ist, wobei in der Basisdarstellung $2c=n\cdot2a$ für jedes Element $2c$ des Ideals die ungeraden $n$ nur Abkürzungen für die endlichen Summen bedeuten.

Die Primideale des Bereichs sind gegeben durch $\ideal P_0=\ideal D=(2)$ und $\ideal P=(2p)$, wo $p$ eine ungerade Primzahl bedeutet; es ist also jedes Primideal durch $\ideal P_0$, aber durch kein weiteres Primideal teilbar. Die primären Ideale sind gegeben durch
\[
 \ideal Q_0=(2\cdot2^{e_0})
 \qquad\text{und}\qquad
 \ideal Q_p=(2p^e);
\]
sie sind zugleich irreduzible Ideale, und nach dem über die Primideale Gesagten sind je zwei zu verschiedenen ungeraden Primzahlen gehörige gegenseitig relativprim, aber kein $\ideal Q$ ist zu irgendeinem $\ideal Q_0$ relativprim.\footnote{In der Tat folgt aus $2b\cdot2p_1^{e_1}\eqzero{2p_2^{e_2}}$ für ungerade $p_1\ne p_2$ stets $2b\eqzero{2p_2^{e_2}}$; aus $2b\cdot2p^e\eqzero{2\cdot2^{e_0}}$ aber nur $2b=2\cdot2^{e_0-1}\pmod{2\cdot2^{e_0}}$.} Die $\ideal Q_0$ sind also die einzigen nicht-isolierten primären Ideale. Der eindeutigen Zerlegung von $a$ in Primzahlpotenzen entspricht die eindeutige Darstellung des Ideals $(2a)$ durch größte primäre, zugleich irreduzible Ideale:
\[
 (2a)=[(2\cdot2^{e_0}),(2p_1^{e_1}),
\ldots,(2p_\mu^{e_\mu})].
\]
Es sind also hier, im Gegensatz zu den Beispielen aus dem Polynombereich, auch die nicht-isolierten größten primären Ideale eindeutig bestimmt. Für $e_0=0$ handelt es sich zugleich um eine Darstellung durch gegenseitig relativprime Ideale, während für $e_0>0$ das Ideal relativprim-irreduzibel ist.

Während also im Bereich aller ganzen Zahlen die vier verschiedenen Zerlegungen zusammenfallen, ist dies hier nur der Fall für die beiden Zerlegungen in größte primäre und irreduzible einerseits, und bei $e_0>0$ für teilerfremd-irreduzible (jedes Ideal ist teilerfremd-irreduzibel, da der Bereich keine Einheit besitzt) und relativprim-irreduzible andererseits, während bei $e_0=0$ die teilerfremd-irreduziblen und relativprim-irreduziblen Zerlegungen voneinander verschieden werden. Zugleich bietet sich hier schon ein Beispiel dafür, daß ein Primideal durch ein anderes teilbar sein kann, ohne damit identisch zu sein; allgemeiner dafür, daß aus der Teilbarkeit nicht die Produktdarstellung folgt. Das letztere - eine Folge davon, daß der Bereich keine Einheit enthält - ist auch der Grund dafür, daß keine eindeutige Produktdarstellung der Zahlen des Bereichs durch irreduzible Zahlen des Bereichs existiert, obwohl jedes Ideal ein Hauptideal wird; die Einführung des kleinsten gemeinsamen Vielfachen erweist sich also hier als notwendig. Es sei noch bemerkt, daß die Verhältnisse genau die gleichen bleiben, wenn statt aller geraden Zahlen alle durch eine feste Primzahl oder Primzahlpotenz teilbaren Zahlen zugrunde gelegt werden.

Dagegen werden auch noch irreduzible und primäre Ideale verschieden, wenn $\Sigma$ aus allen durch eine zusammengesetzte Zahl $g=p_1^{\nu_1}\cdots p_s^{\nu_s}$ teilbaren Zahlen besteht. Es wird wieder jedes Ideal ein Hauptideal $(g\cdot a)$, und die Primideale sind wieder gegeben durch $\ideal P_0=\ideal D=(g)$; $\ideal P=(g\cdot p)$, wo $p$ eine von den in $g$ aufgehenden Primzahlen verschiedene Primzahl bedeutet. Dagegen sind die irreduziblen Ideale gegeben durch $\ideal B_{\lambda_i}=(g\cdot p_i^{\lambda_i})$ und $\ideal B_e=(g\cdot p^e)$; die primären Ideale durch $\ideal Q_e=\ideal B_e$ und durch die von den irreduziblen verschiedenen
\[
 \ideal Q_{\lambda_1\ldots\lambda_s}=(g\cdot p_1^{\lambda_1}\cdots p_s^{\lambda_s}),
\]
wobei die $\ideal B_{\lambda_i}$ und $\ideal Q_{\lambda_1\ldots\lambda_s}$ alle dasselbe zugehörige Primideal $\ideal P_0=(g)$ besitzen. Auch hier besteht Eindeutigkeit der Zerlegung in irreduzible Ideale, und folglich auch für die Zerlegung in größte primäre Ideale; es sind also wieder auch die nicht-isolierten Ideale eindeutig bestimmt.

2. Ein Beispiel eines nicht-kommutativen Ringbereichs bietet die in der Arbeit Noether--Schmeidler behandelte Idealtheorie in nicht-kommutativen Polynombereichen. Es handelt sich insbesondere um ``vollständig-reduzible'' Ideale, d. h. solche, bei denen die Komponenten der Zerlegung paarweise teilerfremd sind und keine echten Teiler besitzen; die Komponenten sind also a fortiori irreduzibel. Es ergibt sich somit in Ergänzung der dort bewiesenen Isomorphie nach \S~9 noch die Anzahlgleichheit der Komponenten bei zwei verschiedenen Zerlegungen. Damit ist für die als Spezialfall der Arbeit sich ergebende Zerlegung der Systeme partieller oder gewöhnlicher linearer Differentialausdrücke ein Resultat gewonnen, das selbst in dem bekannten Fall eines gewöhnlichen linearen Differentialausdrucks nicht bemerkt gewesen zu sein scheint.

Zugleich liefert das System $T$ aller Restklassen eines festen Ideals $M$ zusammen mit dem nicht-kommutativen Polynombereich $\Sigma$ einen Doppelbereich $(\Sigma,T)$, wo $T$ Moduleigenschaft in bezug auf $\Sigma$ hat. Denn die Differenz zweier Restklassen ist wieder eine Restklasse, und ebenso das Produkt einer Restklasse mit einem beliebigen Polynom, während das Produkt zweier Restklassen nicht existiert (a. a. O. \S~3). Die dort als ``Untergruppen'' bezeichneten Systeme von Restklassen bilden somit Beispiele von Moduln in Doppelbereichen $(\Sigma,T)$, wo der zugrunde gelegte Ringbereich $\Sigma$ nicht-kommutativ ist.

\subsection*{\S~12. Beispiel aus der Elementarteilertheorie.}

Es handelt sich um eine durch die allgemeinen Entwicklungen bedingte Auffassung der Elementarteilertheorie, die aber selbst als bekannt vorausgesetzt wird.

Sei $\Sigma$ der Bereich aller ganzzahligen Matrizen aus $n^2$ Elementen, für die Addition und Multiplikation in dem für Matrizen gebräuchlichen Sinn definiert sei. $\Sigma$ wird also ein nicht-kommutativer Ringbereich; die Ideale sind somit im allgemeinen einseitig, die zweiseitigen nur als Spezialfall darin enthalten.\footnote{Die Idealtheorie dieser Bereiche bildet den Gegenstand der Arbeiten von Du Pasquier: Zahlentheorie der Tettarionen, Dissertation Zürich, Vierteljahrsschr. d. Naturf. Ges. Zürich, 51 (1906); Zur Theorie der Tettarionenideale, ibid., 52 (1907). Der Inhalt der zweiten Arbeit ist der Nachweis, daß jedes Ideal ein Hauptideal wird.}

Wir zeigen zuerst, daß jedes Ideal ein Hauptideal wird. Dazu ordnen wir bei rechtsseitigen Idealen jeder Matrix $A=(a_{ik})$ einen Modul
\[
 A=(a_{11}\xi_1+\cdots+a_{1n}\xi_n,
 \ldots,
 a_{n1}\xi_1+\cdots+a_{nn}\xi_n)
\]
aus ganzzahligen Linearformen zu. Umgekehrt entspricht diesem Modul jede Matrix, die eine Basis von $A$ liefert, also neben $A$ auch $UA$, wo $U$ unimodular ist. Allgemeiner entspricht dem Produkt $PA$ ein Modul $B$, der Vielfaches von $A$ wird. Eine einzelne Linearform aus $A$ wird durch solche $P$ gegeben, die nur eine von Null verschiedene Zeile enthalten.

Seien nun $A_1,A_2,
\ldots,A_\nu,
\ldots$ alle Elemente eines Ideals $\ideal M$; $A_1,A_2,
\ldots,A_\nu,
\ldots$ die ihnen zugeordneten Moduln, $A$ deren größter gemeinsamer Teiler und $UA$ die allgemeinste, diesem Modul $A$ zugeordnete Matrix. Jeder einzelnen Linearform aus $A$ entspricht dann nach der Definition des größten gemeinsamen Teilers eine Matrix $P_1A_{i_1}+\cdots+P_\sigma A_{i_\sigma}$, wo nach dem obigen die $P$ nur eine von Null verschiedene Zeile besitzen. Daraus folgt, daß auch die einer Basis von $A$ entsprechende Matrix $A$ eine solche Darstellung, jetzt mit allgemeinem $P$, zuläßt und folglich ein Element aus $\ideal M$ wird. Da ferner jeder Modul $A_i$ durch $A$ teilbar ist, wird jede Matrix $A_i$ durch $A$ teilbar; $A$ und allgemein $UA$ bildet eine Basis von $\ideal M$. Handelt es sich um linksseitige Ideale, so sind entsprechend die Spalten jeder Matrix als Basis eines Moduls zu betrachten; jedes Ideal wird ein Hauptideal, für das neben $A$ auch $AV$ eine Basis wird, unter $V$ eine beliebige unimodulare Matrix verstanden.

Wir legen nun im folgenden zum Zusammenhang mit der Elementarteilertheorie zweiseitige Ideale zugrunde, bei denen also insbesondere neben $A$ auch $PAQ$ zum Ideal gehört. Dann ist die allgemeinste Basis eines solchen Ideals nach dem obigen gegeben durch $UAV$, wo $U$ und $V$ unimodular sind. Die Basiselemente erschöpfen also eine Klasse äquivalenter Matrizen,\footnote{Den Basiselementen der einseitigen Ideale entsprechen Rechts- bzw. Linksklassen.} und es besteht eineindeutige Beziehung zwischen Ideal und Klasse, folglich auch zwischen Ideal und Elementarteilersystem $(a_1\mid a_2\mid\cdots\mid a_n)$ der Klasse, wo die $a_i$ bekanntlich nicht-negative ganze Zahlen sind, von denen jede in der folgenden aufgeht. Die in der durch die Elementarteiler bedingten Normalform auftretende Matrix der Klasse läßt sich somit als spezielle Basis des Ideals auffassen; aus der Teilbarkeit der Ideale bzw. Klassen folgt die der Elementarteiler, und umgekehrt.

Nun hat aber Du Pasquier a. a. O. \S~11 gezeigt, daß bei jedem zweiseitigen Ideal der Rang $n$ ist und alle Elementarteiler übereinstimmen. Um den Fall allgemeiner Elementarteiler betrachten zu können, müssen wir daher nicht von den Idealen, sondern direkt von den zweiseitigen Klassen ausgehen (die ebenfalls durch große deutsche Buchstaben bezeichnet werden sollen). Eine Klasse $\mathfrak A=UAV$ ist dabei also durch eine andere $\mathfrak B$ teilbar, wenn $A=PBQ$ wird.

Allgemein gilt: \emph{Das kleinste gemeinsame Vielfache (der größte gemeinsame Teiler) zweier Klassen wird erhalten durch Bildung des kleinsten gemeinsamen Vielfachen (des größten gemeinsamen Teilers) der entsprechenden Elementarteilersysteme}.\footnote{In der Arbeit: Zur Theorie der Moduln, Math. Ann. 52 (1899), S. 1 definiert E. Steinitz das kleinste gemeinsame Vielfache (den größten gemeinsamen Teiler) von Klassen durch kleinste gemeinsame Vielfache und größte gemeinsame Teiler der Elementarsysteme. Unabhängig davon findet sich das kleinste gemeinsame Vielfache von Klassen als ``Kongruenzkomposition'' bei H. Brandt, Komposition der binären quadratischen Formen relativ einer Grundform, J. f. M. 150 (1919), S. 1.} Denn seien
\[
 (a_1\mid a_2\mid\cdots\mid a_n),\quad
 (b_1\mid b_2\mid\cdots\mid b_n),\quad
 (c_1\mid c_2\mid\cdots\mid c_n),
\]
wo $c_i=[a_i,b_i]$, die Elementarteilersysteme von $\mathfrak A$, $\mathfrak B$, $\mathfrak C^*$, und sei $\mathfrak C=[\mathfrak A,\mathfrak B]$. Dann ist $\mathfrak C^*$ durch $\mathfrak A$ und $\mathfrak B$, also durch $\mathfrak C$ teilbar; umgekehrt sind die Elementarteiler von $\mathfrak C$ durch die von $\mathfrak C^*$ teilbar, also ist $\mathfrak C$ durch $\mathfrak C^*$ teilbar und somit $\mathfrak C=\mathfrak C^*$. Entsprechend ergibt sich der Beweis für den größten gemeinsamen Teiler.

Der eindeutigen Darstellung der Elementarteiler $a_i$ als kleinstes gemeinsames Vielfaches von Primzahlpotenzen entspricht somit eine Darstellung von $\mathfrak A$ als kleinstes gemeinsames Vielfaches von Klassen $\mathfrak Q$, deren Elementarteiler gegeben sind durch die Potenzen einer Primzahl, in Zeichen
\[
 \mathfrak Q\sim(p^{r_1}\mid p^{r_2}\mid\cdots\mid p^{r_\rho}\mid0\mid\cdots\mid0),
 \qquad r_1\le r_2\le\cdots\le r_\rho.
\]
Ist insbesondere der Rang $\rho$ gleich $n$, so hat man hier eine Zerlegung in teilerfremde und teilerfremd-irreduzible Klassen, die trotz des nicht-kommutativen Bereichs eindeutig ist.

Die Klassen $\mathfrak Q$ lassen sich weiter zerlegen in Klassen, die dem Elementarteilersystem entsprechen:
\[
\begin{gathered}
 \mathfrak B_1\sim(p^{r_1}\mid\cdots\mid p^{r_1}),\quad
 \mathfrak B_2\sim(1\mid p^{r_2}\mid\cdots\mid p^{r_2}),\quad \ldots,\\
 \mathfrak B_\nu\sim(1\mid\cdots\mid1\mid p^{r_\nu}\mid\cdots\mid p^{r_\nu}),\quad
 \mathfrak B_{\rho+1}\sim(1\mid\cdots\mid1\mid0\mid\cdots\mid0),
\end{gathered}
\]
wo in $\mathfrak B_\nu$ jeweils $(\nu-1)$-mal die Zahl $1$ auftritt. Ist dabei etwa $r_1=r_2=\cdots=r_\mu=0$, so werden $\mathfrak B_1,
\ldots,\mathfrak B_\mu$ gleich der Einheitsklasse und sind folglich bei der Zerlegung wegzulassen; das gleiche gilt für $\mathfrak B_{\rho+1}$, wenn $\rho=n$ ist. Ist ferner etwa $r_\nu=r_{\nu+1}=\cdots=r_{\nu+\lambda}$, so werden $\mathfrak B_{\nu+1},
\ldots,\mathfrak B_{\nu+\lambda}$ echte Teiler von $\mathfrak B_\nu$, sind also gleichfalls auszuscheiden. Bezeichnen $\mathfrak B_{i_1},
\ldots,\mathfrak B_{i_k}$ die übrigbleibenden, die jetzt eine kürzeste Darstellung ergeben, so wird
\[
 \mathfrak Q=[\mathfrak B_{i_1},
\ldots,\mathfrak B_{i_k}]
\]
die eindeutige Zerlegung von $\mathfrak Q$ in irreduzible Klassen. Sei nämlich
\[
 \mathfrak B_\nu\sim(1\mid\cdots\mid1\mid p^{r_\nu}\mid\cdots\mid p^{r_\nu})
\]
darstellbar als kleinstes gemeinsames Vielfaches von
\[
 \mathfrak C\sim(1\mid\cdots\mid1\mid p^{s_\nu}\mid\cdots\mid p^{s_i})
 \quad\text{und}\quad
 \mathfrak D\sim(1\mid\cdots\mid1\mid p^{t_\nu}\mid\cdots\mid p^{t_i}).
\]
Dann muß im Elementarteilersystem von $\mathfrak C$ und $\mathfrak D$ an den ersten $(\nu-1)$ Stellen die Zahl $1$ stehen; an $\nu$-ter Stelle muß der Exponent $s_\nu$ oder $t_\nu$, etwa $s_\nu$, gleich $r_\nu$ werden. Da aber
\[
 r_\nu=s_\nu\le s_{\nu+1}\le\cdots\le s_i\le r_\nu
\]
wird, so kommt $\mathfrak C=\mathfrak B_\nu$, also ist $\mathfrak B_\nu$ irreduzibel.\footnote{Dagegen sind die $\mathfrak B$ noch in einseitige Klassen reduzibel; hier besteht keine eindeutige Beziehung mehr zwischen Elementarteilern und Klasse, und damit auch keine eindeutige Zerlegung in irreduzible einseitige Klassen. Das zeigt etwa das folgende mir von H. Brandt angegebene Beispiel der Zerlegung in rechtsseitige Klassen (wo die Klassen durch eine Basismatrix dargestellt sind, bzw. durch den entsprechenden Modul):
\[
\begin{gathered}
 \mathfrak B=[\mathfrak C_1,\mathfrak C_2]=[\mathfrak D_1,\mathfrak D_2],\quad
 \mathfrak B\sim\begin{pmatrix}p&0\\0&p\end{pmatrix},\quad
 \mathfrak C_1\sim\begin{pmatrix}1&0\\0&p\end{pmatrix},\quad
 \mathfrak C_2\sim\begin{pmatrix}p&0\\0&1\end{pmatrix},\\
 \mathfrak D_1\sim\begin{pmatrix}1&0\\0&p\end{pmatrix},
 \quad
 \mathfrak D_2\sim\begin{pmatrix}p&0\\(p-1)&1\end{pmatrix}.
\end{gathered}
\]
In der Tat besitzen die Moduln $(\xi,p\eta)$; $(p\xi,\eta)$; $(p\xi,(p-1)\xi+\eta)$ jeweils als echten Teiler nur den Modul $(\xi,\eta)$, sind also irreduzibel und sind ferner voneinander verschieden.} Dasselbe gilt für $\mathfrak B_{\rho+1}$, wo $p$ durch $0$ zu ersetzen ist. Jede dieser irreduziblen Klassen gibt aber, wie die Bildung des kleinsten gemeinsamen Vielfachen zeigt, einen bestimmten Exponenten, und die Stelle, wo dieser Exponent im Elementarteilersystem von $\mathfrak Q$, bzw. $\mathfrak A$ zum erstenmal auftritt, während $\mathfrak B_{\rho+1}$ den Rang angibt. Da diese Zahlen durch das Elementarteilersystem von $\mathfrak Q$, bzw. $\mathfrak A$ eindeutig festgelegt sind, und da die Beziehung zwischen Elementarteilern und Klasse eineindeutig ist, ist somit die Zerlegung von $\mathfrak Q$, und ebenso die einer beliebigen Klasse, in irreduzible Klassen, eindeutig.

Zusammenfassend läßt sich sagen: \emph{Jede zweiseitige Klasse $\mathfrak A$ aus ganzzahligen Matrizen von beschränkter Elementenzahl läßt sich eindeutig darstellen als kleinstes gemeinsames Vielfaches von endlich vielen irreduziblen zweiseitigen Klassen. Jede irreduzible Klasse repräsentiert dabei einen festen Primteiler des Elementarteilersystems von $\mathfrak A$, einen zugehörigen Exponenten und die Stelle, wo dieser Exponent zum erstenmal auftritt. Die dem Teiler $0$ entsprechende irreduzible Klasse gibt den Rang von $\mathfrak A$ an.}

\begin{flushleft}
Erlangen, Oktober 1920.
\end{flushleft}
\begin{center}
(Eingegangen am 16. 10. 1920.)
\end{center}

\end{document}
